\section{Introduction}

Software localization has become a critical aspect of global software development, enabling applications to reach diverse linguistic markets. Machine translation (MT) tools have emerged as valuable assets in this process, offering rapid and cost-effective translation solutions. However, the quality of MT systems varies significantly across different tools and text types, particularly when dealing with domain-specific content such as technical documentation and user interface (UI) strings.

\subsection{Motivation}

Technical texts and UI strings present unique challenges for machine translation systems:

\begin{itemize}
    \item \textbf{Technical Documentation}: Contains domain-specific terminology, complex sentence structures, and requires high accuracy for comprehension.
    \item \textbf{UI Strings}: Often short, context-dependent, and may contain placeholders (e.g., \texttt{\{variable\}}, \texttt{\%s}) that must be preserved.
    \item \textbf{Error Messages}: Critical for user experience and must convey precise meaning.
\end{itemize}

While numerous MT tools are available (Google Translate, DeepL, Microsoft Translator, Amazon Translate), developers and localization teams lack comprehensive comparative data to guide their tool selection, especially for the English-Turkish language pair.

\subsection{Research Questions}

This study addresses the following research questions:

\begin{enumerate}
    \item[\textbf{RQ1:}] Which machine translation tool achieves the highest quality for technical documentation translation?
    \item[\textbf{RQ2:}] Is there a significant performance difference between tools for UI string translation?
    \item[\textbf{RQ3:}] How does text length affect translation quality across different tools?
    \item[\textbf{RQ4:}] Which tool maintains better technical terminology consistency?
    \item[\textbf{RQ5:}] How well do different tools preserve placeholders and special characters?
    \item[\textbf{RQ6:}] What is the cost-quality trade-off for each tool?
\end{enumerate}

\subsection{Contributions}

This paper makes the following contributions:

\begin{itemize}
    \item A curated dataset of 50,000+ English-Turkish parallel segments from software localization sources, categorized by text type (technical documentation, UI strings, error messages).
    \item Comprehensive evaluation of four major MT tools using multiple automatic metrics (BLEU, METEOR, TER, chrF++).
    \item Statistical analysis demonstrating significant performance differences across tools and text categories.
    \item An open-source web application for interactive MT tool comparison.
    \item Practical recommendations for software developers and localization teams.
\end{itemize}

\subsection{Paper Organization}

The remainder of this paper is organized as follows: Section II reviews related work on MT evaluation and software localization. Section III describes our methodology, including dataset construction and evaluation metrics. Section IV details the dataset characteristics. Section V presents our experimental setup and results. Section VI discusses findings and implications. Section VII concludes with future work directions.
